% Drawing 147/03-8 -- folder 147, batch 03, archive page 54 (PDF page 55).
% Transcription: les trois graphes de (66), dessinés l'un sous l'autre : une boucle fermée sur un sommet marqué « 1 » ; deux sommets marqués « 1 » reliés par deux arcs ; un cercle avec une queue, le bout de la queue marqué « 2 ».
% Kind: three MD-graphs with first Betti number 1. Method 2 (schematic redraw from the facsimile). Status: ok.
% Notes: Three graphs, one under the other, as on the page. (1) A loop (a drop) closed on a vertex marked 1. (2) Two vertices, each marked 1, joined by two arcs. (3) A circle with a tail: a segment from a free end marked 2 to a vertex on the circle. In each graph the edges are drawn between the same computed vertices. These are the three graphs of the table numbered (65) on p. 53 (drawing 03-5), drawn again; the transcription sets the table in words. Overlaid 2026-11-12 (Opus 5.5) on a 400 dpi colour render: the loop of the first graph was about 3 units too flat and the circle of the third about 3 units too high; both traced again through points of the page. Redrawn 2026-11-12 (issue 128): to check again.
% Coordinates: 1 unit = 1 mm = 1/100 of the crop width (bbox in page fractions: [0.29, 0.315, 0.47, 0.465]).
% Redrawn by Claude (claude-fable-5-2), issue #128, 2026-11-12; not checked by a human; not compiled (no TeX in the session).
% Overlaid by Claude (claude-opus-5-5) 2026-11-12 (issue #128): rasterised from its own coordinates over the colour facsimile and checked or corrected. Not compiled. Status: ok.
% Rights: reproduces part of the work -- behind RENDER_TRANSCRIPTS, never published by itself.
\documentclass[tikz,border=3pt]{standalone}
\usepackage{amsmath,amssymb}
\usetikzlibrary{patterns,arrows.meta}
\tikzset{pencil/.style={black!45}, struck/.style={black!35}, mpt/.style={circle,draw,inner sep=1.3pt}, pt/.style={circle,fill,inner sep=0.9pt}, opt/.style={circle,draw,fill=white,inner sep=1.1pt}}
\begin{document}
\begin{tikzpicture}[x=1mm,y=1mm,line cap=round,line join=round,font=\scriptsize]
\coordinate (A) at (28,-20);
\draw plot[smooth, tension=0.6] coordinates {(28,-20) (36,-12.5) (46,-10) (54,-11.5) (58.5,-17) (59,-22) (56,-28.5) (48,-31) (38,-29) (28,-20)};
\node[pt] at (A) {};
\node[inner sep=1pt] at (23.5,-13) {\small $1$};
\coordinate (B) at (26,-59);
\coordinate (C) at (60,-59);
\draw (B) to[out=55,in=125] (C);
\draw (B) to[out=-55,in=-125] (C);
\node[pt] at (B) {};
\node[pt] at (C) {};
\node[inner sep=1pt] at (19,-53.5) {\small $1$};
\node[inner sep=1pt] at (68,-56) {\small $1$};
\coordinate (D) at (21,-95.5);
\coordinate (E) at (50.5,-95.5);
\draw (D) -- (E);
\draw (67.5,-97) ellipse [x radius=17, y radius=15.2];
\node[pt] at (D) {};
\node[pt] at (E) {};
\node[inner sep=1pt] at (20,-88) {\small $2$};
\end{tikzpicture}
\end{document}
